\section{From possible samples to a sampling distribution}

The previous section showed that the complete sample
\[
(X_1,X_2,\ldots,X_n)
\]
is a random object. It has many possible values, and the sampling model assigns
probabilities to those values.

A statistic is obtained by applying the same rule to every possible sample. For
example, the sample mean is the rule
\[
(x_1,\ldots,x_n)
\longmapsto
\frac{x_1+\cdots+x_n}{n}.
\]
Before the sample is observed, this rule produces the random variable
\[
\overline X=\frac{X_1+\cdots+X_n}{n}.
\]
Because different samples may give different values of \(\overline X\), the
sample mean has its own probability distribution.

\begin{definitionbox}
The \textbf{sampling distribution} of a statistic is the probability law of that
statistic under repeated sampling from a specified model.
\end{definitionbox}

The word ``sampling'' is important. The distribution does not describe the
values inside one observed dataset. It describes the values that the statistic
could take over all possible samples generated by the model.

\subsection*{A complete example with samples of size two}

Return to two independent Bernoulli observations with
\[
P(X_i=1)=p,
\qquad
P(X_i=0)=1-p.
\]
The possible ordered samples are
\[
(0,0),\qquad(0,1),\qquad(1,0),\qquad(1,1).
\]
Now define the sample mean
\[
\overline X=\frac{X_1+X_2}{2}.
\]
For zero--one data, this mean is also the proportion of successes in the sample.
Applying the rule to every possible sample gives
\[
\begin{array}{c|c|c}
\text{possible sample} & \text{probability} & \overline X \\
\hline
(0,0) & (1-p)^2 & 0 \\
(0,1) & (1-p)p & \frac12 \\
(1,0) & p(1-p) & \frac12 \\
(1,1) & p^2 & 1
\end{array}
\]

The statistic has only three possible values:
\[
0,\qquad \frac12,\qquad 1.
\]
The middle value is produced by two different samples. Therefore
\[
P(\overline X=0)=(1-p)^2,
\]
\[
P\left(\overline X=\frac12\right)
=(1-p)p+p(1-p)=2p(1-p),
\]
and
\[
P(\overline X=1)=p^2.
\]
These probabilities add to one:
\[
(1-p)^2+2p(1-p)+p^2=1.
\]
This is the sampling distribution of \(\overline X\).

\begin{examplebox}
For \(p=0.6\), the sampling distribution is
\[
P(\overline X=0)=0.16,
\]
\[
P\left(\overline X=\frac12\right)=0.48,
\]
\[
P(\overline X=1)=0.36.
\]
The value \(1/2\) is most likely because it can be produced by either \((0,1)\)
or \((1,0)\).
\end{examplebox}

\subsection*{Several samples may produce the same statistic}

A statistic usually records less information than the complete sample. The
ordered samples
\[
(0,1)
\qquad\text{and}\qquad
(1,0)
\]
are different, but both have sample mean \(1/2\). The statistic groups together
all samples that produce the same summary value.

This is the same modelling idea that appeared earlier in the course. A detailed
outcome is transformed into a coarser description:
\[
\text{complete sample}
\longrightarrow
\text{statistic value}.
\]
The probability of a statistic value is obtained by adding the probabilities of
all samples that lead to that value.

\begin{remarkbox}
A sampling distribution is not created by listing possible numerical values and
assigning equal probabilities to them. Its probabilities are inherited from the
probabilities of the samples that produce those values.
\end{remarkbox}

\subsection*{The same sample can produce many statistics}

From one random sample we may compute many different summaries. For example,
from
\[
(X_1,X_2,X_3)
\]
we may form
\[
\overline X=\frac{X_1+X_2+X_3}{3},
\]
\[
M=\operatorname{median}(X_1,X_2,X_3),
\]
or
\[
R=\max(X_1,X_2,X_3)-\min(X_1,X_2,X_3).
\]
Each rule produces a different random variable and therefore a different
sampling distribution.

The purpose of this chapter is not to derive every such distribution. The
important idea is more general:
\[
\text{random sample}
\longrightarrow
\text{statistic}
\longrightarrow
\text{sampling distribution}.
\]

\begin{summarybox}
A statistic has a probability law because it is a function of a random sample.
To construct its sampling distribution in a small model:
\begin{enumerate}
    \item list the possible samples;
    \item assign probabilities to those samples;
    \item compute the statistic for each sample;
    \item add the probabilities of samples giving the same statistic value.
\end{enumerate}
\end{summarybox}
